\documentclass[12pt]{article}

\usepackage[margin=2cm]{geometry}
\usepackage[utf8]{inputenc}
\usepackage{amsfonts}
\usepackage{physics}
\usepackage{amsmath}
\usepackage{amssymb}
\usepackage{mathtools}
\usepackage{csquotes}
\usepackage[backend=bibtex, sorting=none, style=numeric-comp]{biblatex}
\addbibresource{references.bib}
\usepackage[english]{babel}
\usepackage{hyperref}

\newcommand{\qp}{q\,p}
\newcommand{\vphi}{\varphi}

\newcommand{\U}{\mathcal{U}} % propagator
\newcommand{\PF}{\mathcal{P}} % perron-frobenius
\newcommand{\kr}{\kappa} % krylov state
\newcommand{\M}{\mathcal{M}} % map
\newcommand{\Ms}[1][1]{\mathcal{M}^{#1}} % power of the map

\newcommand{\hrho}{\hat{\rho}} % hatted rho
\newcommand{\hkr}{\hat{\kappa}} % hatted krylov state

\newcommand{\sket}[1]{|#1)}
\newcommand{\sbra}[1]{(#1|}
\newcommand{\sbraket}[2]{\left(#1|#2\right)}
\newcommand{\smel}[3]{\left(#1|#2|#3\right)}

\newcommand{\bnt}[2]{\beta_{#1\,#2}}
\newcommand{\atn}[2]{\alpha_{#1\,#2}}

\begin{document}

This document contains definitions and analytic derivations that are used in the code base but remain otherwise unexplained in the manuscript,
with the aim of making the functionality of the code more transparent.

\section{Jacobi theta and ``controlled" Jacobi theta}
In the quantization of classical maps on the torus the Jacobi theta function
$$
\Theta(v|\tau) = \sum_{\nu\in\mathbb{Z}} e^{i\pi\tau \nu^2 + 2iv \nu}
$$
appears frequently.
For instance, a coherent state on the torus $\ket{z}$ is expressed in the position basis $\{\ket{q_n}\}_{n=0}^{N-1}$ as \cite{PLeboeuf_1990}
\begin{equation}
    \braket{q_n}{z} = A e^{-\pi N (q_n^2 + z^2 - 2\sqrt{2} z q_n)}
    \Theta\left(i\pi N (q_n - \sqrt{2}z) | i N\right),
    \label{eq:qn_z}
\end{equation}
where $A$ is a normalization constant and $N$ the dimension of the Hilbert space.
Let $z = (q + ip)/\sqrt{2}$ (so $\ket{z}$ is a coherent state centred at a point $(q,p)$ in phase space),
writing the Jacobi theta explicitly one gets
\begin{equation}
    \Theta_n(q, p) \equiv \Theta\left(i\pi N (q_n - \sqrt{2}z) | i N\right) =
    e^{\pi N (q - q_n)^2} \sum_{\nu\in\mathbb{Z}} e^{-\pi N (\nu - q + q_n)^2} e^{i 2\pi N p \nu}.
    \label{eq:theta}
\end{equation}
In practice, even for rather modest values of $N$ the sum is very well approximated by only a few terms.
In the code we use $-2 \leq \nu \leq 2$.
This expression is badly behaved because of the exponential factor outside the sum which causes overflows in the numerics,
so instead of using \eqref{eq:theta} we define the ``controlled" Jacobi theta
\begin{equation}
\boxed{
    \tilde{\Theta}_n(q, p) =
    \sum_{\nu\in\mathbb{Z}} e^{-\pi N (\nu - q + q_n)^2} e^{i 2\pi N p \nu}
}
    \label{eq:controlled_theta}
\end{equation}
and absorb the exponential factor into the rest of \eqref{eq:qn_z},
which after some rearranging yields
$$
    \braket{q_n}{\qp} = \tilde{A} e^{-i\pi N p (q - 2 q_n)} \tilde{\Theta}_n(q, p)
$$
where we have denoted $\ket{z} \equiv \ket{\qp}$.
It is then clear that after normalizing we have
\begin{equation}
    \boxed{
    \braket{q_n}{\qp} = e^{-i\pi N p (q - 2 q_n)}
    \frac{\tilde{\Theta}_n(q, p)}{\sqrt{\sum_{k=0}^{N-1} \abs{\tilde{\Theta}_k(q, p)}^2 }}
}
    \label{eq:controlled_qn_z},
\end{equation}
which is the expression used in the code.


\section{Density matrix of a ``coherent ensemble"}
We refer to as a ``coherent ensemble" the ``diagonal" representation of a density matrix in the coherent basis
\begin{equation}
    \hrho = \iint dqdp\, f(q,p) \ketbra{\qp},
    \label{eq:coherent_ensemble}
\end{equation}
where $\ket{\qp}$ is the coherent state centred at a point $(q, p)$ in phase space.

\subsection{Harmonic oscillator}
In the excitation number representation this is:
$$
\rho_{nm} = \mel{n}{\hrho}{m} = \iint dqdp\, f(q,p) \braket{n}{\qp}\braket{\qp}{m}.
$$
One has
$$
\braket{n}{\qp} = \sqrt{ P_n\left(\frac{q^2 + p^2}{2\hbar}\right) } e^{in\vphi},
$$
where $P_n(x)= \frac{1}{n!} x^n e^{-x}$ is a Poisson probability mass function with mean $x$
and $\varphi = \arctan(p/q)$ is the argument of the point $(q, p)$ in phase space.
After changing into polar coordinates:
$$
\rho_{nm} = \iint drd\vphi\, r f(r\cos(\vphi), r\sin(\vphi))
\sqrt{ P_n\left(\frac{r^2}{2\hbar}\right) P_m\left(\frac{r^2}{2\hbar}\right) } e^{i\vphi(n-m)}.
$$
Finally, letting $u = r/\sqrt{2\hbar}$ and expanding:
$$
\rho_{nm} = \frac{2\hbar}{\sqrt{n! m!}} \iint dud\vphi\,
f\left(\sqrt{2\hbar}u\cos(\vphi), \sqrt{2\hbar}u\sin(\vphi)\right)
u^{n+m+1} e^{-u^2} e^{i\vphi(n - m)}.
$$
This is the integral we compute numerically, but the previous expression is not well behaved because of the factorials and powers of $u$.
To avoid overflows we rewrite it as
$$
\boxed{
\rho_{nm} = 2\hbar \iint dud\vphi\,
f\left(\sqrt{2\hbar}u\cos(\vphi), \sqrt{2\hbar}u\sin(\vphi)\right) e^{L_{nm}(u)} e^{i\vphi(n - m)}
},
$$
where
$$
L_{nm}(u) = -u^2 + (n+m+1)\ln(u) - \frac{1}{2}[\ln(\Gamma(n+1))+\ln(\Gamma(m+1)))],
$$
with $\Gamma(n+1) = n!$ the Gamma function.
The logarithm of the Gamma function is implemented in SciPy.

\subsection{Harper map}
In the position basis we have
$$
\rho_{nm} = \iint dqdp\, f(q,p) \braket{q_n}{\qp}\braket{\qp}{q_m}.
$$
Here $\ket{\qp}$ refers to a coherent state defined on the torus.
Using Eq.~\eqref{eq:controlled_qn_z}, the integral we compute numerically is
$$
\boxed{
\rho_{nm} = \iint dqdp\, f(q,p) e^{i 2\pi N p (q_n - q_m)}
\frac{\tilde{\Theta}_n(q, p) \tilde{\Theta}_m(q, p)^*}{\sum_{k=0}^{N-1} \abs{\tilde{\Theta}_k(q, p)}^2 }
},
$$
where $a^*$ denotes $a$'s complex conjugation.

\section{Propagator in the Krylov basis from the wavefunction}
Given the quantum propagator $\U$, one can usually compute its matrix representation in the Krylov basis $\{\hkr_n\}$ by explicitly evaluating $\smel{\hkr_n}{\U}{\hkr_m}$.
However, in a classical setting the propagator is the Perron-Frobenius operator $\PF$ which acts on a distribution $\rho$ over phase space as $\PF \rho(q, p) = \rho(\M^{-1}(q, p))$.
To evaluate this numerically one needs an explicit functional expression for $\rho$.
In practice, we do these evaluations to get the evolution of the distribution $\{\rho_t(q, p)\}$,
and using the Gram-Schmidt procedure we get the associated Krylov distributions $\{\kr_n(q, p)\}$.
We can't compute $\PF \kr_n(q, p)$ because in practice we don't have an explicit functional expression for $\kr_n$,
but an array of data.
Thus, we can't evaluate $\smel{\kr_n}{\PF}{\kr_m}$ explicitly.
However, with the data we are able to calculate the Krylov wavefunction $\bnt{n}{t}$ and this in turn allows us to compute $\smel{\kr_n}{\PF}{\kr_m}$ indirectly as follows.

Recall that the Krylov states in terms of the evolution are
\begin{equation}
    \kr_n = \sum_{t=0}^n \atn{t}{n}\,\rho_t,
    \label{eq:kn_rhot}
\end{equation}
and conversely
\begin{equation}
    \rho_t = \sum_{n=0}^t \bnt{n}{t}\,\kr_{n}
    \label{eq:rhot_kn}.
\end{equation}
Thus
$$
\bnt{n}{t} = \sbraket{\kr_n}{\rho_t} = \iint dqdp\, \kr_n(q, p) \rho_t(q, p),
$$
which we can compute numerically with the data.
Inserting Eq.\eqref{eq:kn_rhot} in the right Krylov state in $\smel{\kr_n}{\PF}{\kr_m}$
and using that $\PF \sket{\rho_t} = \sket{\rho_{t+1}}$ we have
$$
\boxed{
\smel{\kr_n}{\PF}{\kr_m} = \sum_{t=0}^m \bnt{n}{t+1} \atn{t}{m}
}.
$$
Now we need $\atn{t}{n}$.
Using the orthogonality of the Krylov states $\sbraket{\kr_n}{\kr_m} = \delta_{nm}$ and Eqs.\eqref{eq:kn_rhot} and \eqref{eq:rhot_kn} it is easy to see that
$$
\sum_{t=0}^m \bnt{n}{t} \atn{t}{m} = \delta_{nm},
$$
which means that $\beta$ and $\alpha$ are recursively inverse matrices,
i.e., let $\beta^{(m)}$ be the $m\times m$ matrix with elements $(\beta^{(m)})_{n t} = \bnt{n}{t}$ and $\alpha^{(m)}$ the $m\times m$ matrix with elements $(\alpha^{(m)})_{tm} = \atn{t}{m}$,
then $\beta^{(m)} \alpha^{(m)} = \mathbb{I}_{m\times m}$.
That is to say, we can calculate $\atn{t}{n}$ simply by inverting the matrix $\beta$.

% \clearpage
% \pagebreak
\printbibliography

\end{document}
